% ============================== CHAPTER 1 ============================
\setlecture{Lecture I}\renewcommand{\tabname}{I}
\chapter{Functions of Several Variables}

In first-year calculus a function eats one number and returns one number. In real
problems almost nothing depends on a single quantity. The temperature in a room depends
on three coordinates $(x,y,z)$, the price of a phone depends on its storage, its camera
and its brand, the loss of a neural network depends on millions of weights. This chapter
sets up the language for such functions and builds the first tools to study them:
pictures (graphs and level curves), partial derivatives, the gradient, and the precise
meaning of \emph{differentiable}.

\section{Functions on sets of points}

\mnote{A point of $\R^n$ is a list $X=(x_1,\dots,x_n)$. The same list can be read as a
point (a position) or as a vector (an arrow from the origin). We use both readings.}

\begin{definition}{Functions on sets of points}
Let $S$ be a set of points in $n$-space. A \emph{function defined on $S$} associates to
each point of $S$ a number:
\[
 f : S\subset \R^n \longrightarrow \R, \qquad X\longmapsto f(X).
\]
The set $S$ is called the \emph{domain of definition} of $f$.
\end{definition}

In practice the domain is often not written: it is understood to be the largest set of
points where the formula makes sense. Finding it means listing every operation in the
formula and asking where it is allowed.

\begin{recipe}
To find the natural domain of a formula, check each of these:
\begin{itemize}
  \item \textbf{Division}: the denominator must not be $0$.
  \item \textbf{Even roots} $\sqrt{\;\cdot\;}$: the inside must be $\ge 0$.
  \item \textbf{Logarithms} $\ln(\cdot)$: the inside must be $>0$.
  \item \textbf{$\arcsin$, $\arccos$}: the inside must be in $[-1,1]$.
\end{itemize}
Then intersect all the conditions and \emph{draw} the resulting region.
\end{recipe}

\begin{example}{distance from the origin}
Let $f(x,y)=x^2+y^2$. Every point of the plane is allowed, so $\Dom(f)=\R^2$.
Geometrically,
\[ f(x,y)=\norm{(x,y)}^2, \]
the square of the distance from $(x,y)$ to the origin. For instance $f(3,4)=25$, and
indeed $(3,4)$ is at distance $5$ from $O$.
\end{example}

\begin{example}{a function not defined at the origin}
Let
\[ f(x,y)=\frac{x^2-y^2}{x^2+y^2}. \]
The denominator vanishes only when $x=y=0$. Hence the origin is removed:
\[ \Dom(f)=\R^2\setminus\{(0,0)\}. \]
We will see in Section~\ref{sec:level} that this innocent-looking function behaves
very strangely near the missing point.
\end{example}

\begin{example}{three variables}
In $3$-space, $f(x,y,z)=x^2-\sin(xyz)+yz^3$ is defined everywhere. Writing
$X=(x,y,z)$ we simply write $f(X)$. For example
$f(1,0,2)=1-\sin 0+0=1$.
\end{example}

\begin{example}{finding a domain}
Find the domain of $f(x,y)=\sqrt{4-x^2-y^2}$.

\textbf{Step 1.} The only restriction comes from the square root: we need
$4-x^2-y^2\ge 0$.

\textbf{Step 2.} Rearranging, $x^2+y^2\le 4$, so
\[ \Dom(f)=\{(x,y)\in\R^2 : x^2+y^2\le 4\}. \]

\textbf{Step 3.} This is the \emph{closed disc} of radius $2$ centred at the origin. Its
boundary is the circle $x^2+y^2=4$, and the boundary is included because of the
``$\ge$''.

\begin{center}
\begin{tikzpicture}
\begin{axis}[bp, width=6.2cm, height=6.2cm, xmin=-3, xmax=3, ymin=-3, ymax=3,
   xtick={-2,2}, ytick={-2,2}, xlabel=$x$, ylabel=$y$, axis equal image]
  \addplot[domain=0:360, samples=90, draw=plum, fill=plum, fill opacity=0.18]
     ({2*cos(x)},{2*sin(x)}) -- cycle;
  \node[font=\small, text=plum] at (axis cs:0.9,0.8) {$x^2+y^2\le 4$};
\end{axis}
\end{tikzpicture}
\end{center}
\end{example}

\begin{example}{a half-plane and a hole}
Find the domain of $g(x,y)=\dfrac{\ln(x-y)}{\sqrt{x}}$.

The logarithm needs $x-y>0$, that is $y<x$. The root in the denominator needs $x\ge 0$
and, being in a denominator, $x\neq 0$; together, $x>0$. So
\[ \Dom(g)=\{(x,y): x>0,\ y<x\}, \]
the part of the right half-plane lying strictly below the line $y=x$. Neither boundary
line is included.
\end{example}

\begin{reallife}{cost functions}
A factory's weekly cost might be $C(L,K)=40L+90K+\sqrt{LK}$, where $L$ is hours of labour
and $K$ hours of machine time. The formula makes sense for all $L,K\ge 0$, but the
\emph{physical} domain is smaller: a week has at most $168$ hours of machine time. Real
problems often restrict the domain further than the formula does.
\end{reallife}

\section{Graphs of functions}

\begin{definition}{Graph}
For a function of $n$ variables $f(x_1,\dots,x_n)$, the \emph{graph} of $f$ is the set of
points in $(n+1)$-space of the form
\[ \big(x_1,\dots,x_n,\,f(x_1,\dots,x_n)\big). \]
\end{definition}

\begin{itemize}
  \item $n=1$: the graph is a curve in the plane, points $(x,f(x))$.
  \item $n=2$: the graph is a \emph{surface} in $3$-space, points $(x,y,f(x,y))$. Think
  of $f(x,y)$ as the height of a landscape above the point $(x,y)$ on the floor.
  \item $n=3$: the graph lives in $4$-space, which we cannot draw. We need another tool:
  level sets.
\end{itemize}

\begin{center}
\begin{tikzpicture}
\begin{axis}[bp3d, width=10.5cm, height=9cm, xlabel=$x$, ylabel=$y$, zlabel=$z$,
   xmin=-2.2, xmax=2.2, ymin=-2.2, ymax=2.2, zmin=0, zmax=8.5]
  \addplot3[plumsurf, domain=0:2, y domain=0:360, samples=16, samples y=36,
     z buffer=sort] ({x*cos(y)},{x*sin(y)},{x^2});
  \addplot3[only marks, mark=*, coral, mark size=2.2pt] coordinates {(1,1,0)};
  \addplot3[coral, dashed, line width=0.9pt] coordinates {(1,1,0) (1,1,2)};
  \addplot3[only marks, mark=*, coral, mark size=2.2pt] coordinates {(1,1,2)};
  \node[font=\scriptsize, text=coraldark, anchor=north] at (axis cs:1,1,0) {$(x,y)$};
  \node[font=\scriptsize, text=coraldark, anchor=west] at (axis cs:1.05,1,2.1) {$(x,y,f(x,y))$};
\end{axis}
\end{tikzpicture}

{\small\color{navy} The graph of $f(x,y)=x^2+y^2$ is a paraboloid: above each point
$(x,y)$ of the floor sits the point at height $f(x,y)$.}
\end{center}

\section{Level curves and level surfaces}\label{sec:level}

\begin{intuition}
A hiking map is flat, yet it shows mountains. It does so with \emph{contour lines}:
each line joins points at the same altitude. A level curve is exactly a contour line of
the landscape $z=f(x,y)$. You slice the surface horizontally at height $c$ and look
straight down at the cut.
\end{intuition}

\begin{definition}{Level curves}
For a function $f(x,y)$ and a number $c$, the \emph{level curve of $f$ at level $c$} is
the set of points $(x,y)$ where $f$ takes the value $c$:
\[ \{(x,y) : f(x,y)=c\}. \]
\end{definition}

\begin{itemize}
  \item A graph records the height above the plane.
  \item A level curve records where the height is constant.
  \item Drawing several level curves together gives a \emph{contour map}.
\end{itemize}

\subsection{The paraboloid $f(x,y)=x^2+y^2$}

\begin{example}{level curves of $x^2+y^2$}
The level curves are $x^2+y^2=c$.
\begin{itemize}
  \item If $c<0$: no point satisfies it (a sum of squares is never negative). The level
  set is empty.
  \item If $c=0$: only the origin.
  \item If $c>0$: the circle of radius $\sqrt c$.
\end{itemize}
Notice that the radii $1,\sqrt2,\sqrt3,2$ for $c=1,2,3,4$ get \emph{closer together} as
$c$ grows: the paraboloid becomes steeper as we move away from the origin.
\end{example}

\begin{center}
\begin{tikzpicture}
\begin{axis}[bp, width=5.7cm, height=5.7cm, xmin=-2.6, xmax=2.6, ymin=-2.6, ymax=2.6,
   xtick={-2,-1,1,2}, ytick={-2,-1,1,2}, xlabel=$x$, ylabel=$y$, axis equal image]
  \foreach \c [evaluate=\c as \r using sqrt(\c)] in {1,2,3,4} {
    \edef\tmp{\noexpand\addplot[domain=0:360, samples=80,
       draw=plum!\the\numexpr 40+15*\c\relax!coral] ({\r*cos(x)},{\r*sin(x)});}\tmp }
  \node[font=\scriptsize, text=plum, fill=white, inner sep=1pt] at (axis cs:0.72,0.72) {$c{=}1$};
  \node[font=\scriptsize, text=plum, fill=white, inner sep=1pt] at (axis cs:1.42,1.42) {$c{=}4$};
\end{axis}
\end{tikzpicture}\hspace{0.3cm}
\begin{tikzpicture}
\begin{axis}[bp3d, width=7.4cm, height=7cm, xlabel=$x$, ylabel=$y$, zlabel=$z$,
   xmin=-2.3, xmax=2.3, ymin=-2.3, ymax=2.3, zmin=0, zmax=5.2]
  \addplot3[plumsurf, opacity=0.55, fill opacity=0.55, domain=0:2.1, y domain=0:360,
    samples=12, samples y=32, z buffer=sort] ({x*cos(y)},{x*sin(y)},{x^2});
  \foreach \c [evaluate=\c as \r using sqrt(\c)] in {1,2,3,4} {
    \edef\tmp{\noexpand\addplot3[cyanbp, line width=1.1pt, domain=0:360, samples=60, samples y=1]
       ({\r*cos(x)},{\r*sin(x)},\c);
     \noexpand\addplot3[cyanbp!60, line width=0.6pt, dashed, domain=0:360, samples=60, samples y=1]
       ({\r*cos(x)},{\r*sin(x)},0);}\tmp }
\end{axis}
\end{tikzpicture}

{\small\color{navy} Left: the contour map. Right: the same circles as horizontal slices
of the surface (solid) and their shadows on the floor (dashed).}
\end{center}

\subsection{The saddle $f(x,y)=x^2-y^2$}

\begin{example}{level curves of $x^2-y^2$}
The level curves are $x^2-y^2=c$.
\begin{itemize}
  \item $c>0$: hyperbolas opening left and right (they cross the $x$-axis at $\pm\sqrt c$).
  \item $c<0$: hyperbolas opening up and down (they cross the $y$-axis at $\pm\sqrt{-c}$).
  \item $c=0$: $x^2=y^2$, that is the two lines $y=x$ and $y=-x$.
\end{itemize}
The graph is a \emph{saddle} (a horse saddle, or a Pringles chip): going along the
$x$-axis you climb, going along the $y$-axis you descend.
\end{example}

\begin{center}
\begin{tikzpicture}
\begin{axis}[bp, width=5.7cm, height=5.7cm, xmin=-2.5, xmax=2.5, ymin=-2.5, ymax=2.5,
   xtick={-2,-1,1,2}, ytick={-2,-1,1,2}, xlabel=$x$, ylabel=$y$, axis equal image]
  \foreach \c/\T in {0.5/1.75, 1.5/1.15, 3/0.62} {
    \edef\tmp{%
    \noexpand\addplot[plum, domain=-\T:\T, samples=40] ({sqrt(\c)*cosh(x)},{sqrt(\c)*sinh(x)});
    \noexpand\addplot[plum, domain=-\T:\T, samples=40] ({-sqrt(\c)*cosh(x)},{sqrt(\c)*sinh(x)});
    \noexpand\addplot[coral, domain=-\T:\T, samples=40] ({sqrt(\c)*sinh(x)},{sqrt(\c)*cosh(x)});
    \noexpand\addplot[coral, domain=-\T:\T, samples=40] ({sqrt(\c)*sinh(x)},{-sqrt(\c)*cosh(x)});}\tmp }
  \addplot[navy, dashed, domain=-2.5:2.5] {x};
  \addplot[navy, dashed, domain=-2.5:2.5] {-x};
  \node[font=\scriptsize, text=plum, fill=white, inner sep=1pt] at (axis cs:2.05,0.25) {$c>0$};
  \node[font=\scriptsize, text=coraldark, fill=white, inner sep=1pt] at (axis cs:0.25,2.1) {$c<0$};
\end{axis}
\end{tikzpicture}\hspace{0.3cm}
\begin{tikzpicture}
\begin{axis}[bp3d, width=7.4cm, height=7cm, xlabel=$x$, ylabel=$y$, zlabel=$z$,
   xmin=-1.8, xmax=1.8, ymin=-1.8, ymax=1.8, zmin=-3, zmax=3, view={-30}{24}]
  \addplot3[plumsurf, domain=-1.5:1.5, y domain=-1.5:1.5, samples=22] {x^2-y^2};
\end{axis}
\end{tikzpicture}
\end{center}

\subsection{A function whose level curves are lines}

\begin{example}{level curves of $\frac{x^2-y^2}{x^2+y^2}$}
Let $f(x,y)=\dfrac{x^2-y^2}{x^2+y^2}$ for $(x,y)\neq(0,0)$.

\textbf{Step 1: write the level equation.} $f(x,y)=c$ means
$x^2-y^2=c(x^2+y^2)$.

\textbf{Step 2: collect the $x$ and $y$ terms.}
\[ x^2-cx^2 = y^2+cy^2 \quad\Longleftrightarrow\quad x^2(1-c)=y^2(1+c). \]

\textbf{Step 3: solve.} For $x\neq 0$ and $c\neq -1$,
\[ \frac{y^2}{x^2}=\frac{1-c}{1+c}, \qquad\text{so}\qquad
   y=\pm a x,\quad a=\sqrt{\frac{1-c}{1+c}} . \]
This requires $\frac{1-c}{1+c}\ge0$, that is $-1<c\le 1$.

\textbf{Step 4: special values.}
\begin{itemize}
  \item $c=1$: $a=0$, the level curve is the $x$-axis (without $O$).
  \item $c=0$: $a=1$, the lines $y=\pm x$.
  \item $c=-1$: the equation becomes $2x^2=0$, so $x=0$: the $y$-axis (without $O$).
  \item $\abs c>1$: empty. Indeed $\abs{x^2-y^2}\le x^2+y^2$, so $-1\le f\le 1$ always.
\end{itemize}
So \emph{every} level curve is a pair of lines through the origin, with the origin
removed. All the values between $-1$ and $1$ crowd into every neighbourhood of $O$. This
is the first sign that $f$ cannot be extended continuously to the origin.
\end{example}

\begin{center}
\begin{tikzpicture}
\begin{axis}[bp, width=6.3cm, height=6.3cm, xmin=-2.2, xmax=2.2, ymin=-2.2, ymax=2.2,
   xtick={-1,1}, ytick={-1,1}, xlabel=$x$, ylabel=$y$, axis equal image]
  \addplot[coral, domain=-2.2:2.2, line width=1.6pt] {0};
  \addplot[coral!70!plum, domain=-2.2:2.2] {0.4142*x};
  \addplot[coral!70!plum, domain=-2.2:2.2] {-0.4142*x};
  \addplot[plummid, domain=-2.2:2.2] {x};
  \addplot[plummid, domain=-2.2:2.2] {-x};
  \addplot[plum, domain=-0.9:0.9] {2.414*x};
  \addplot[plum, domain=-0.9:0.9] {-2.414*x};
  \addplot[plumdark, line width=1.6pt] coordinates {(0,-2.2) (0,2.2)};
  \node[odot] at (axis cs:0,0) {};
  \node[font=\scriptsize, fill=white, inner sep=1pt, text=coraldark] at (axis cs:1.9,0.22) {$c{=}1$};
  \node[font=\scriptsize, fill=white, inner sep=1pt, text=plummid] at (axis cs:1.75,1.45) {$c{=}0$};
  \node[font=\scriptsize, fill=white, inner sep=1pt, text=plumdark] at (axis cs:0.42,2.0) {$c{=}{-}1$};
\end{axis}
\end{tikzpicture}

{\small\color{navy} The rays from the origin are level curves: $c=1$ on the $x$-axis,
$c=\tfrac{1}{\sqrt2}$ on $y=\pm(\sqrt2-1)x$, $c=0$ on the diagonals, $c=-1$ on the $y$-axis.}
\end{center}

\subsection{Reading a contour map}

\begin{center}
\begin{tikzpicture}[scale=0.9]
  \draw[bpgrid, step=0.5] (-3,-2.3) grid (3,2.3);
  \foreach \s/\c in {2.6/40, 1.9/60, 1.35/80, 0.9/100} {
    \draw[plum!\c!coral, line width=1.1pt, xscale=1.25]
      plot[smooth cycle, tension=0.8] coordinates {(-\s,0) (-0.6*\s,0.75*\s) (0.5*\s,0.7*\s) (\s*0.95,0.1*\s) (0.45*\s,-0.65*\s) (-0.55*\s,-0.62*\s)};}
  \node[font=\scriptsize, text=plum] at (-3.0,0.25) {1};
  \node[font=\scriptsize, text=plum] at (-2.15,0.25) {2};
  \node[font=\scriptsize, text=plum] at (-1.45,0.25) {3};
  \node[font=\scriptsize, text=plum] at (-0.85,0.25) {4};
  \draw[vecc] (1.0,-0.1) -- (2.9,-0.1) node[right, font=\scriptsize]{gentle};
  \draw[vec] (0.0,0.75) -- (0,1.95) node[above, font=\scriptsize]{steep};
\end{tikzpicture}
\end{center}

Suppose these are level curves of $f$ for the values $1,2,3,4$.
\begin{itemize}
  \item Moving inward, $f$ increases: this is a hill with its top inside the smallest curve.
  \item Each curve contains points with the same value.
  \item \textbf{Where is the graph steepest?} Where the curves are \emph{closest together}.
  Between two consecutive curves $f$ changes by $1$; if they are close, this change
  happens over a short distance, so the slope is large.
\end{itemize}

\begin{reallife}{weather maps and topography}
On a weather map the level curves of pressure are called \emph{isobars}. Where isobars
are packed together the pressure changes fast, and the wind is strong. On a hiking map,
packed contour lines warn you of a cliff. The same reading applies to any contour plot,
including the loss landscapes plotted in machine learning.
\end{reallife}

\subsection{Level surfaces}

For a function of three variables we cannot draw the graph, but we can draw where
the function is constant.

\begin{definition}{Level surface}
For $f(x,y,z)$ and a number $c$, the \emph{level surface} at level $c$ is the set
$\{(x,y,z): f(x,y,z)=c\}$.
\end{definition}

\mfig{\begin{tikzpicture}[scale=0.62]
  \foreach \r/\o in {1/1, 1.6/0.7, 2.2/0.45} {
   \shade[ball color=plum!40!white, opacity=\o] (0,0) circle (\r);
   \draw[plum, line width=0.5pt] (0,0) circle (\r);
   \draw[plum, line width=0.4pt, dashed] (-\r,0) arc (180:0:\r cm and 0.3*\r cm);
   \draw[plum, line width=0.5pt] (-\r,0) arc (180:360:\r cm and 0.3*\r cm);}
  \node[font=\scriptsize, text=navy, align=center] at (0,-2.75) {$x^2+y^2+z^2=c$\\nested spheres};
\end{tikzpicture}}

For $f(x,y,z)=x^2+y^2+z^2$ the level surfaces are spheres of radius $\sqrt c$ centred at
the origin, nested like the layers of an onion.

\begin{reallife}{equipotentials and isotherms}
If $f$ is an electric potential, its level surfaces are the \emph{equipotential surfaces};
moving a charge along one of them costs no work. If $f$ is temperature, the level
surfaces are \emph{isothermal surfaces}. In both cases we describe the field by the
places where it takes a fixed value.
\end{reallife}

\section{Open sets: room around every point}

\begin{intuition}
To differentiate $f$ at $x$ in one variable you compute $f(x+h)$ for small $h$ on both
sides. So $f$ must be defined on a little interval around $x$. In several variables
you want to move a little in \emph{every} direction, so you need a little disc (or ball)
around the point. Sets where every point has such a ball are called \emph{open}.
\end{intuition}

\begin{definition}{Balls and spheres}
Let $P\in\R^n$ and $a>0$.
\[
\text{open ball: } \norm{X-P}<a,\qquad
\text{closed ball: } \norm{X-P}\le a,
\]
\[ \text{sphere: } \norm{X-P}=a. \]
\end{definition}

\begin{itemize}
  \item In $\R$: an open ball is an open interval $(P-a,P+a)$.
  \item In $\R^2$: an open ball is an open disc and the sphere is a circle.
  \item In $\R^3$: the words mean what they usually mean.
\end{itemize}

\textbf{Scaling.} If $S_1$ is the unit sphere centred at $O$ and $X\in S_1$, then for
$a>0$, $\norm{aX}=a\norm X=a$. Multiplying by $a$ stretches the unit sphere onto the sphere of
radius $a$. The same holds for balls.

\begin{definition}{Open set}
A set $U\subset\R^n$ is \emph{open} if for every $P\in U$ there exists $a>0$ such that
the open ball centred at $P$ of radius $a$ is contained in $U$.
\end{definition}

\begin{center}
\begin{tikzpicture}[scale=0.8]
  \begin{scope}
   \draw[bpgrid, step=0.5] (-2.6,-2.6) grid (2.6,2.6);
   \fill[plum, opacity=0.15] (0,0) circle (2);
   \draw[plum, line width=1.2pt, dashed] (0,0) circle (2);
   \node[pdot] (p) at (1.25,0.8) {};
   \fill[coral, opacity=0.3] (p) circle (0.5);
   \draw[coral, line width=0.9pt] (p) circle (0.5);
   \node[font=\scriptsize, text=plumdark, above left] at (p) {$P$};
   \node[font=\small, text=navy] at (0,-3.0) {open disc $x^2+y^2<4$};
   \node[font=\scriptsize, text=navy] at (0,-3.5) {every $P$ has room around it};
  \end{scope}
  \begin{scope}[xshift=6.5cm]
   \draw[bpgrid, step=0.5] (-2.6,-2.6) grid (2.6,2.6);
   \fill[plum, opacity=0.15] (0,0) circle (2);
   \draw[plum, line width=1.4pt] (0,0) circle (2);
   \node[pdot] (p) at (45:2) {};
   \fill[coral, opacity=0.3] (p) circle (0.5);
   \draw[coral, line width=0.9pt] (p) circle (0.5);
   \node[font=\scriptsize, text=plumdark, above right] at (p) {$P$};
   \node[font=\small, text=navy] at (0,-3.0) {closed disc $x^2+y^2\le 4$};
   \node[font=\scriptsize, text=navy] at (0,-3.5) {at a boundary point every ball escapes};
  \end{scope}
\end{tikzpicture}
\end{center}

\begin{example}{which sets are open}
\begin{itemize}
  \item The first quadrant without the axes, $\{x>0,\ y>0\}$, is open in $\R^2$: around
  $(x,y)$ take the ball of radius $\min(x,y)$.
  \item The $x$-axis is \emph{not} open as a subset of $\R^2$: any disc around a point of
  the axis contains points with $y\neq0$.
  \item The same axis, viewed as the whole space $\R$, \emph{is} open.
  \item The interval $(-1,1)$ is open in $\R$, but not as a subset of $\R^2$.
\end{itemize}
\end{example}

\begin{trap}
``Open'' is not a property of a set alone; it depends on the \emph{ambient space}. Always
ask: open as a subset of which $\R^n$? Also, ``not open'' does not mean ``closed'':
a half-open interval $[0,1)$ is neither.
\end{trap}

\begin{theorem}{An open ball is open}
Let $a>0$. The open ball $U=\{X:\norm X<a\}$ is an open set.
\end{theorem}

\begin{intuition}
If $P$ is at distance $b$ from the centre, it is at distance $a-b$ from the edge.
A ball of radius $a-b$ around $P$ therefore cannot reach the edge. The triangle
inequality turns this picture into a proof.
\begin{center}
\begin{tikzpicture}[scale=0.55]
  \draw[plum, dashed, line width=0.9pt] (0,0) circle (2.5);
  \fill[plum, opacity=0.08] (0,0) circle (2.5);
  \node[pdot] (o) at (0,0) {};
  \node[dot] (p) at (1.4,0.6) {};
  \draw[navy, line width=0.6pt] (o) -- node[below, font=\tiny]{$b$} (p);
  \fill[coral, opacity=0.25] (p) circle ({2.5-1.523});
  \draw[coral] (p) circle ({2.5-1.523});
  \draw[coral, line width=0.6pt] (p) -- ++(23.2:{2.5-1.523}) node[midway, above, font=\tiny]{$c$};
  \node[font=\tiny, below left] at (o) {$O$};
  \node[font=\tiny, above] at (p) {$P$};
  \node[font=\scriptsize, text=navy] at (0,-3.0) {$c=a-b$};
\end{tikzpicture}
\end{center}
\end{intuition}



\begin{proof}
Take $P\in U$ and put $b=\norm P$; then $b<a$. Choose $c=a-b>0$. Let $X$ be any point
with $\norm{X-P}<c$. By the triangle inequality,
\[
 \norm X=\norm{(X-P)+P}\le\norm{X-P}+\norm P<(a-b)+b=a .
\]
So $X\in U$. Thus the ball of radius $c$ around $P$ lies in $U$, and $U$ is open.
\end{proof}

The point of all this: if $X\in U$ and $U$ is open, then the nearby points $X+H$ with
$\norm H$ small are also in $U$. In particular we can change \emph{one} coordinate a
little and stay in $U$. That is exactly what partial derivatives need.

\section{Partial derivatives}

\begin{intuition}
Standing on a hillside at $(x_0,y_0)$, you can ask: if I walk due east, how steep is it?
If I walk due north, how steep? Each question freezes one coordinate and moves the
other, so it is an ordinary one-variable derivative. These two slopes are the partial
derivatives $\partial f/\partial x$ and $\partial f/\partial y$.
\end{intuition}

\begin{definition}{The $i$-th partial derivative}
Let $f$ be defined on an open set $U\subset\R^n$ and $X=(x_1,\dots,x_n)\in U$. If the
limit exists, the $i$-th partial derivative of $f$ at $X$ is
\[
 D_if(X)=\pd{f}{x_i}(X)=\lim_{h\to0}\frac{f(x_1,\dots,x_i+h,\dots,x_n)-f(x_1,\dots,x_n)}{h}.
\]
\end{definition}

\begin{itemize}
  \item Only the variable $x_i$ changes; all the others are kept fixed.
  \item For two variables:
  \[
  \pd fx=\lim_{h\to0}\frac{f(x+h,y)-f(x,y)}{h},\qquad
  \pd fy=\lim_{k\to0}\frac{f(x,y+k)-f(x,y)}{k}.
  \]
\end{itemize}

\subsection{Partial derivatives as slices}

To compute $\pd fx(x_0,y_0)$: freeze $y=y_0$, get the one-variable function
$\varphi(x)=f(x,y_0)$, and differentiate: $\varphi'(x_0)=\pd fx(x_0,y_0)$. Similarly
$\psi(y)=f(x_0,y)$ gives $\psi'(y_0)=\pd fy(x_0,y_0)$.

\begin{center}
\begin{tikzpicture}
\begin{axis}[bp3d, width=11cm, height=9.2cm, xlabel=$x$, ylabel=$y$, zlabel=$z$,
   xmin=-0.2, xmax=2.3, ymin=-0.2, ymax=2.3, zmin=0, zmax=6.5, view={-40}{25}]
  \addplot3[plumsurf, opacity=0.6, fill opacity=0.6, domain=0:2, y domain=0:2, samples=16]
     {4-0.5*(x-1)^2-0.75*(y-1)^2+0.6*x};
  % slice y = 1
  \addplot3[cyanbp, line width=1.8pt, domain=0:2, samples=40, samples y=1]
     (x,1,{4-0.5*(x-1)^2+0.6*x});
  % slice x = 1
  \addplot3[coral, line width=1.8pt, domain=0:2, samples=40, samples y=1]
     (1,x,{4.6-0.75*(x-1)^2});
  \addplot3[only marks, mark=*, mark size=2.6pt, white, mark options={draw=plumdark, line width=1pt}]
     coordinates {(1,1,4.6)};
  % tangent along x : slope 0.6
  \addplot3[cyanbp, dashed, line width=1pt, domain=0.4:1.6, samples=2, samples y=1]
     (x,1,{4.6+0.6*(x-1)});
  % footprints
  \addplot3[cyanbp!50, line width=0.6pt, dashed, domain=0:2, samples=2, samples y=1] (x,1,0);
  \addplot3[coral!50, line width=0.6pt, dashed, domain=0:2, samples=2, samples y=1] (1,x,0);
  \node[font=\scriptsize, text=cyanbp, anchor=west] at (axis cs:2,1,5.1) {$y=y_0$ fixed};
  \node[font=\scriptsize, text=coraldark, anchor=south] at (axis cs:1,2.1,4.0) {$x=x_0$ fixed};
\end{axis}
\end{tikzpicture}

{\small\color{navy} The blue curve is the slice $y=y_0$; its slope at the marked point
is $\partial f/\partial x$ (dashed tangent). The coral curve is the slice $x=x_0$; its
slope gives $\partial f/\partial y$.}
\end{center}

\begin{example}{computing partials}
Let $f(x,y)=x^2y^3$.
\begin{itemize}
  \item For $\pd fx$, treat $y$ as a constant: $y^3$ is a constant factor, so
  $\pd fx=2x\,y^3$.
  \item For $\pd fy$, treat $x$ as a constant: $\pd fy=3x^2y^2$.
\end{itemize}
The partial derivatives are themselves functions of $(x,y)$.
\end{example}

\begin{trap}
$D_if(P)$ means: first compute the \emph{function} $D_if$, then evaluate it at $P$.
It does \emph{not} mean $D_i\big(f(P)\big)$: $f(P)$ is a number and its derivative is $0$.
For example with $f=x^2y^3$ and $P=(1,2)$: $D_1f(P)=2\cdot1\cdot 8=16$, while
``$D_1(f(1,2))$'' $=D_1(8)=0$.
\end{trap}

\begin{example}{evaluating a partial derivative}
Let $f(x,y)=\sin(xy)$. Find $D_2f(1,\pi)$ and $D_2f(3,\pi/4)$.

\textbf{Step 1.} Differentiate in $y$, with $x$ fixed (chain rule, inner function $xy$,
whose $y$-derivative is $x$):
\[ D_2f(x,y)=x\cos(xy). \]
\textbf{Step 2.} Evaluate:
\[
D_2f(1,\pi)=1\cdot\cos\pi=-1,\qquad
D_2f\!\left(3,\tfrac\pi4\right)=3\cos\tfrac{3\pi}{4}=-\tfrac{3}{\sqrt2}.
\]
\end{example}

\begin{example}{a quotient}
Let $f(x,y)=\dfrac{x}{x^2+y^2}$. With $y$ fixed, the quotient rule gives
\begin{align*}
\pd fx&=\frac{(x^2+y^2)\cdot1-x\cdot2x}{(x^2+y^2)^2}=\frac{y^2-x^2}{(x^2+y^2)^2},\\
\pd fy&=x\cdot\frac{-2y}{(x^2+y^2)^2}=\frac{-2xy}{(x^2+y^2)^2}.
\end{align*}
\end{example}

\begin{example}{using the limit definition}
Let $f(x,y)=\sqrt[3]{x^3+y^3}$. Find $D_1f(0,0)$.

The formula obtained by differentiating, $\frac{x^2}{(x^3+y^3)^{2/3}}$, is undefined at
the origin, so we go back to the definition:
\[
D_1f(0,0)=\lim_{h\to0}\frac{f(h,0)-f(0,0)}{h}=\lim_{h\to0}\frac{\sqrt[3]{h^3}}{h}
=\lim_{h\to0}\frac hh=1 .
\]
Lesson: when the derivative formula breaks down at a point, use the definition.
\end{example}

\section{The gradient}

\begin{definition}{Gradient}
Assume the partial derivatives of $f$ exist at $X$. The \emph{gradient} of $f$ at $X$ is
the vector of partial derivatives:
\[
\operatorname{grad} f(X)=\grad f(X)=\big(D_1f(X),\dots,D_nf(X)\big).
\]
For two variables $\grad f(x,y)=\left(\pd fx,\pd fy\right)$; for three,
$\grad f=\left(\pd fx,\pd fy,\pd fz\right)$.
\end{definition}

The gradient assigns a \emph{vector} to each point. We read $\grad f(X)$ as
$(\grad f)(X)$: compute the gradient function, then evaluate.

\begin{example}{a gradient}
For $f(x,y)=x^2y^3$, $\grad f(x,y)=(2xy^3,\;3x^2y^2)$, so
$\grad f(1,2)=(2\cdot8,\;3\cdot4)=(16,12)$.
\end{example}

\begin{example}{a first picture}
For $f(x,y)=x^2+y^2$, $\grad f(x,y)=(2x,2y)=2(x,y)$. At every point the gradient points
straight away from the origin, and its length $2\norm{(x,y)}$ grows as we go out.
Compare with the level curves: the gradient is \emph{perpendicular} to them and
points towards higher values.
\end{example}

\begin{center}
\begin{tikzpicture}
\begin{axis}[bp, width=6.8cm, height=6.8cm, xmin=-2.4, xmax=2.4, ymin=-2.4, ymax=2.4,
   xtick={-2,-1,1,2}, ytick={-2,-1,1,2}, xlabel=$x$, ylabel=$y$, axis equal image]
  \foreach \r in {0.7,1.4,2.1}
     {\edef\tmp{\noexpand\addplot[plum!70, domain=0:360, samples=70] ({\r*cos(x)},{\r*sin(x)});}\tmp}
  \foreach \xx in {-1.8,-1.2,...,1.9} { \foreach \yy in {-1.8,-1.2,...,1.9} {
     \edef\tmp{\noexpand\draw[coral, -{Stealth[length=1.6mm]}, line width=0.9pt]
       (axis cs:\xx,\yy) -- (axis cs:{\xx+0.3*\xx},{\yy+0.3*\yy});}\tmp }}
\end{axis}
\end{tikzpicture}
\end{center}

In Chapter~\ref{ch:chain} we will prove that this is no accident: $\grad f$ always points
in the direction of fastest increase and is always perpendicular to level curves.

\begin{theorem}{Linearity of the gradient}
Let $f,g$ be defined on an open set $U$ with all partial derivatives existing on $U$.
For any number $c$,
\[
\grad(f+g)=\grad f+\grad g,\qquad \grad(cf)=c\,\grad f .
\]
\end{theorem}

\begin{proof}
Each component of the gradient is a partial derivative, which is an ordinary
one-variable derivative with the other variables frozen. One-variable derivatives
satisfy $(f+g)'=f'+g'$ and $(cf)'=cf'$. So for every $i$,
$D_i(f+g)=D_if+D_ig$ and $D_i(cf)=cD_if$. Putting the components together gives the
two vector identities.
\end{proof}

\begin{extra}
The same componentwise argument gives a product rule:
$\grad(fg)=g\,\grad f+f\,\grad g$. For example with $f=x$, $g=y^2$:
$\grad(xy^2)=(y^2,2xy)=y^2(1,0)+x(0,2y)$. This is used constantly in exercises.
\end{extra}

\section{Differentiability}

\subsection{What replaces the Newton quotient?}

Let $f$ be defined on an open set $U$ and $X\in U$. For small $H\neq O$, the point
$X+H$ is in $U$. But the expression
\[ \frac{f(X+H)-f(X)}{H} \]
is meaningless: \textbf{we cannot divide by a vector}. We need a way to say
``$f$ has a derivative'' with no division at all.

\subsection{The one-variable reformulation}

\begin{intuition}
In one variable, being differentiable means the graph looks more and more like a
straight line when you zoom in. Algebraically: the change $f(x+h)-f(x)$ equals a
linear term $f'(x)h$ plus an error that is small \emph{even compared with $h$}. This
statement contains no division, so it can be copied to several variables.
\end{intuition}

In one variable, $f'(x)=\lim_{h\to0}\frac{f(x+h)-f(x)}{h}$. Equivalently,
\[
 f(x+h)-f(x)=f'(x)\,h+\abs h\,g(h),\qquad g(h)\to0 .
\]
\begin{itemize}
  \item Main term: the linear approximation $f'(x)h$.
  \item Error term: $\abs h g(h)$, which is $o(\abs h)$ (smaller than any multiple of $\abs h$).
\end{itemize}

\begin{proposition}{The two forms agree}
If $f(x+h)-f(x)=ah+\abs h g(h)$ with $g(h)\to0$, then $f$ is differentiable at $x$ and
$f'(x)=a$. Conversely, if $f'(x)$ exists, such a $g$ exists with $a=f'(x)$.
\end{proposition}

\begin{proof}
Suppose the formula holds. For $h\neq0$,
\[
\frac{f(x+h)-f(x)}{h}=a+\frac{\abs h}{h}\,g(h).
\]
Since $\abs{\abs h/h}=1$ and $g(h)\to0$, the last term tends to $0$, so the quotient tends
to $a$. Conversely, if $f'(x)$ exists, define for $h\neq0$
\[
g(h)=\frac{1}{\abs h}\big(f(x+h)-f(x)-f'(x)h\big)=\frac{h}{\abs h}\left(\frac{f(x+h)-f(x)}{h}-f'(x)\right),
\]
and $g(0)=0$. The bracket tends to $0$ and $h/\abs h$ has absolute value $1$, so $g(h)\to0$.
\end{proof}

\begin{keyidea}
Differentiability means: \textbf{linear part plus an error of smaller order}.
\end{keyidea}

\subsection{The definition in several variables}

Write $X=(x_1,\dots,x_n)$ and $H=(h_1,\dots,h_n)$, so $X+H=(x_1+h_1,\dots,x_n+h_n)$. In
two variables, $X=(x,y)$, $H=(h,k)$ and $X+H=(x+h,y+k)$.
\begin{itemize}
  \item \emph{Continuity} asks whether $f(X+H)-f(X)\to0$ as $H\to O$.
  \item \emph{Differentiability} asks for a good linear approximation of this difference.
\end{itemize}

\begin{definition}{Differentiability}
$f$ is \emph{differentiable at $X$} if the partial derivatives exist and there is a
function $g(H)$ with $g(H)\to0$ as $H\to O$ such that
\[
 f(X+H)-f(X)=D_1f(X)h_1+\dots+D_nf(X)h_n+\norm H g(H).
\]
Using the gradient:
\[
 \boxed{\,f(X+H)-f(X)=\grad f(X)\cdot H+\norm H\,g(H)\,}
\]
$f$ is \emph{differentiable on $U$} if it is differentiable at every point of $U$.
\end{definition}

In two variables this reads
\[
 f(x+h,y+k)-f(x,y)=\pd fx h+\pd fy k+\norm H g(H),
\]
and the linear part is $\pd fx h+\pd fy k=\grad f(x,y)\cdot(h,k)$.

\begin{center}
\begin{tikzpicture}
\begin{axis}[bp3d, width=10.5cm, height=8.6cm, xlabel=$x$, ylabel=$y$, zlabel=$z$,
   xmin=-0.3, xmax=2.4, ymin=-0.3, ymax=2.4, zmin=0, zmax=8.5, view={-50}{22}]
  \addplot3[plumsurf, opacity=0.75, fill opacity=0.75, domain=0:2, y domain=0:2, samples=16] {x^2+y^2};
  \addplot3[surf, shader=flat, draw=cyanbp, fill=cyanbp, fill opacity=0.3, draw opacity=0.7,
     domain=0.45:1.55, y domain=0.45:1.55, samples=4] {2*x+2*y-2};
  \addplot3[only marks, mark=*, mark size=2.5pt, coral] coordinates {(1,1,2)};
  \node[font=\scriptsize, text=cyanbp, anchor=west] at (axis cs:1.8,0.2,3.0) {linear part};
\end{axis}
\end{tikzpicture}

{\small\color{navy} Near $P=(1,1)$ the surface $z=x^2+y^2$ is almost the plane
$z=2+2(x-1)+2(y-1)$. Differentiability says the gap is $o(\norm H)$.}
\end{center}

\begin{trap}
Read $\grad f(X)$ as $(\grad f)(X)$, never as $\grad\big(f(X)\big)$. The second is the
gradient of a constant, which is $O$.
\end{trap}

\begin{example}{a differentiable function}
Show that $f(x,y)=x^2+y^2$ is differentiable everywhere.

\textbf{Step 1: expand the increment.} With $X=(x,y)$, $H=(h,k)$:
\[
f(X+H)-f(X)=(x+h)^2+(y+k)^2-x^2-y^2=2xh+2yk+h^2+k^2 .
\]
\textbf{Step 2: spot the linear part.} $\grad f(X)=(2x,2y)$, so
$2xh+2yk=\grad f(X)\cdot H$. What is left is $h^2+k^2=\norm H^2$:
\[
f(X+H)-f(X)=\underbrace{\grad f(X)\cdot H}_{\text{linear part}}+\underbrace{\norm H^2}_{\text{error}} .
\]
\textbf{Step 3: check the error is $o(\norm H)$.} Write $\norm H^2=\norm H g(H)$ with
$g(H)=\norm H\to0$. Hence $f$ is differentiable at every point.
\end{example}

\begin{example}{a product}
Show that $f(x,y)=xy$ is differentiable everywhere.

$f(x+h,y+k)-f(x,y)=(x+h)(y+k)-xy=yh+xk+hk$. Since $\grad f=(y,x)$, the linear part is
$yh+xk$ and the error is $hk$. Now $\abs{hk}\le\norm H\cdot\norm H$, so
$g(H)=hk/\norm H$ (with $g(O)=0$) satisfies $\abs{g(H)}\le\norm H\to0$.
\end{example}

\subsection{Partial derivatives inside differentiability}

Look in the direction of the $i$-th coordinate vector $E_i$, that is $H=hE_i=(0,\dots,h,\dots,0)$.
Then $\grad f(X)\cdot H=D_if(X)h$ and $\norm H=\abs h$, so the definition becomes
\[ f(X+hE_i)=f(X)+\pd f{x_i}(X)\,h+\abs h\,g(hE_i). \]
\begin{itemize}
  \item Along coordinate directions, differentiability gives back the partial derivatives.
  \item Partial derivatives are \emph{necessary} for differentiability.
  \item They are \emph{not sufficient} without extra hypotheses (Section~\ref{sec:counter}).
\end{itemize}

\subsection{Tools for estimating remainders}

With $H=(h,k)$, $\norm H=\sqrt{h^2+k^2}$, so
\[ \abs h\le\norm H,\qquad \abs k\le\norm H . \]
Combined with the triangle inequality $\abs{a+b}\le\abs a+\abs b$, this controls most
error terms. For instance
\[ \abs{h^2+hk}\le\abs h^2+\abs h\abs k\le 2\norm H^2 , \]
so $h^2+hk=\norm H g(H)$ with $\abs{g(H)}\le 2\norm H\to0$.

\begin{proposition}{A useful remainder pattern}
If $g_1(H)\to0$ and $g_2(H)\to0$ as $H\to O$, then
$hg_1(H)+kg_2(H)=\norm H g(H)$ for some $g(H)\to 0$.
\end{proposition}

\begin{proof}
For $H\neq O$ factor out $\norm H$:
\[
hg_1(H)+kg_2(H)=\norm H\left(\frac{h}{\norm H}g_1(H)+\frac{k}{\norm H}g_2(H)\right).
\]
Since $\abs{h/\norm H}\le1$ and $\abs{k/\norm H}\le1$, the bracket has absolute value at most
$\abs{g_1(H)}+\abs{g_2(H)}\to0$. Call the bracket $g(H)$.
\end{proof}

\subsection{A sufficient condition}

\begin{theorem}[label=thm:C1]{$C^1$ implies differentiable}
Let $f$ be defined on an open set $U$. Assume that all first partial derivatives exist on
$U$ and are continuous. Then $f$ is differentiable on $U$.
\end{theorem}

\begin{intuition}
To go from $(x,y)$ to $(x+h,y+k)$, walk in an L-shape: first horizontally, then
vertically. Each leg changes only one variable, so the one-variable mean value theorem
turns each change into ``partial derivative times step''. The partial derivatives are
evaluated at slightly wrong points, but continuity says that the mistake tends to $0$.
That mistake is exactly the remainder.
\begin{center}
\begin{tikzpicture}[scale=0.8]
  \draw[bpgrid, step=0.4] (-0.3,-0.3) grid (3.1,2.5);
  \node[pdot, label={[font=\tiny]below:$(x,y)$}] (a) at (0,0) {};
  \node[pdot, label={[font=\tiny]above:$(x,y{+}k)$}] (b) at (0,2) {};
  \node[dot, label={[font=\tiny]above:$(x{+}h,y{+}k)$}] (c) at (2.6,2) {};
  \draw[vecc] (a) -- node[left, font=\tiny]{2nd} (b);
  \draw[vec] (b) -- node[below, font=\tiny]{1st} (c);
  \node[font=\tiny, text=navy] at (1.3,-0.6) {split the increment};
\end{tikzpicture}
\end{center}
\end{intuition}



\begin{proof}
We work in two variables; the general case is the same with more terms. Fix $(x,y)\in U$
and let $H=(h,k)$ be small enough that the whole rectangle with corners $(x,y)$ and
$(x+h,y+k)$ lies in $U$.

\textbf{Step 1: split the increment.} Add and subtract $f(x,y+k)$:
\begin{align*}
f(x+h,y+k)-f(x,y)&=\big[f(x+h,y+k)-f(x,y+k)\big]\\&\quad+\big[f(x,y+k)-f(x,y)\big].
\end{align*}
In the first bracket only $x$ changes ($y+k$ is fixed); in the second only $y$ changes.

\textbf{Step 2: mean value theorem.} Apply the one-variable mean value theorem to
$t\mapsto f(t,y+k)$ on $[x,x+h]$ and to $t\mapsto f(x,t)$ on $[y,y+k]$. There are
numbers $s$ between $x$ and $x+h$, and $t$ between $y$ and $y+k$, with
\begin{align*}
f(x+h,y+k)-f(x,y+k)&=D_1f(s,y+k)\,h,\\ f(x,y+k)-f(x,y)&=D_2f(x,t)\,k .
\end{align*}

\textbf{Step 3: use continuity.} As $H\to O$, $(s,y+k)\to(x,y)$ and $(x,t)\to(x,y)$.
Since $D_1f$ and $D_2f$ are continuous,
\[
D_1f(s,y+k)=D_1f(x,y)+g_1(H),\qquad D_2f(x,t)=D_2f(x,y)+g_2(H),
\]
where $g_1(H)\to0$ and $g_2(H)\to0$.

\textbf{Step 4: collect.} Substituting,
\[
f(x+h,y+k)-f(x,y)=D_1f(x,y)h+D_2f(x,y)k+hg_1(H)+kg_2(H).
\]
The first two terms are $\grad f(x,y)\cdot(h,k)$. By the remainder pattern above, the
last two terms equal $\norm Hg(H)$ with $g(H)\to0$. Therefore
$f(X+H)-f(X)=\grad f(X)\cdot H+\norm Hg(H)$: $f$ is differentiable at $X$.
\end{proof}

\begin{extra}
\textbf{$n$ variables.} Write $f(X+H)-f(X)$ as a telescoping sum, changing one coordinate
at a time:
\begin{align*}
 &\big[f(x_1+h_1,x_2+h_2,\dots,x_n+h_n)-f(x_1,x_2+h_2,\dots,x_n+h_n)\big]\\
 +\,&\big[f(x_1,x_2+h_2,\dots,x_n+h_n)-f(x_1,x_2,x_3+h_3,\dots,x_n+h_n)\big]\\
 +\,&\cdots+\big[f(x_1,\dots,x_{n-1},x_n+h_n)-f(x_1,\dots,x_n)\big].
\end{align*}
Apply the mean value theorem to each bracket and sum the estimates.

\textbf{In practice.} Polynomials, $\sin$, $\cos$, $\exp$, and their sums, products,
quotients (where the denominator is nonzero) and compositions have continuous partial
derivatives, so they are differentiable on their domains. This theorem is how you justify
differentiability in an exam in one line.
\end{extra}

\subsection{Warning: partial derivatives alone are not enough}\label{sec:counter}

\begin{example}{the classic counterexample}
Let
\[
f(x,y)=\frac{xy}{x^2+y^2}\ \text{ if }(x,y)\neq(0,0),\qquad f(0,0)=0 .
\]
\textbf{The partials exist at the origin.} On the $x$-axis $f(h,0)=0$, so
\[ D_1f(0,0)=\lim_{h\to0}\frac{f(h,0)-f(0,0)}{h}=\lim_{h\to0}\frac{0-0}{h}=0, \]
and similarly $D_2f(0,0)=0$.

\textbf{But $f$ is not continuous at the origin.} Approach along lines:
\begin{align*}
\text{along }y=x&:\quad f(x,x)=\frac{x^2}{2x^2}=\frac12,\\
\text{along }y=-x&:\quad f(x,-x)=\frac{-x^2}{2x^2}=-\frac12 .
\end{align*}
Two different limits, so $\lim_{(x,y)\to(0,0)}f(x,y)$ does not exist. Since a differentiable
function is continuous (box below), $f$ is not differentiable at $O$.
\end{example}

\begin{center}
\begin{tikzpicture}
\begin{axis}[bp3d, width=10cm, height=8.4cm, xlabel=$x$, ylabel=$y$, zlabel=$z$,
   xmin=-1.3, xmax=1.3, ymin=-1.3, ymax=1.3, zmin=-0.8, zmax=0.8, view={-30}{30}]
  \addplot3[plumsurf, domain=0.04:1, y domain=0:360, samples=10, samples y=48, z buffer=sort]
    ({x*cos(y)},{x*sin(y)},{0.5*sin(2*y)});
  \addplot3[cyanbp, line width=1.6pt, samples=2, samples y=1, domain=0:1] ({x*0.7071},{x*0.7071},0.5);
  \addplot3[coral, line width=1.6pt, samples=2, samples y=1, domain=0:1] ({x*0.7071},{-x*0.7071},-0.5);
\end{axis}
\end{tikzpicture}

{\small\color{navy} In polar form $f=\tfrac12\sin 2\theta$: the height depends only on the
angle. Along $y=x$ the surface sits at $\tfrac12$ (blue), along $y=-x$ at $-\tfrac12$
(coral), all the way to the origin.}
\end{center}

\begin{extra}
\textbf{Differentiable implies continuous.} If $f(X+H)-f(X)=\grad f(X)\cdot H+\norm H g(H)$,
then by Cauchy-Schwarz $\abs{\grad f(X)\cdot H}\le\norm{\grad f(X)}\norm H\to0$, and
$\norm H g(H)\to0$. So $f(X+H)\to f(X)$.

\medskip
\textbf{The hierarchy to remember:}
\begin{center}
\begin{tikzpicture}[node distance=0.5cm and 0.9cm,
   b/.style={draw=plum, line width=0.8pt, fill=plumtint, font=\small, inner sep=5pt, align=center}]
 \node[b] (c1) {partials exist\\ and are continuous};
 \node[b, right=of c1] (d) {differentiable};
 \node[b, right=of d, yshift=0.6cm] (c) {continuous};
 \node[b, right=of d, yshift=-0.6cm] (p) {partials exist};
 \draw[vec] (c1) -- (d);
 \draw[vec] (d) -- (c.west);
 \draw[vec] (d) -- (p.west);
\end{tikzpicture}
\end{center}
None of the arrows can be reversed. The counterexample above shows that ``partials exist''
does not even give ``continuous''.
\end{extra}

\begin{reallife}{why engineers care}
Linear approximation is how every simulation code works. To predict how the stress in a
beam changes when two dimensions are changed by small amounts $h$ and $k$, an engineer
uses $\Delta\sigma\approx\pd\sigma xh+\pd\sigma yk$. Differentiability is the guarantee that this
prediction gets better and better as the changes get smaller, in every direction at once.
\end{reallife}

\section{Summary}

\begin{keyidea}
\begin{itemize}
  \item A function of several variables maps points of $\R^n$ to numbers. Its natural
  domain is where the formula makes sense.
  \item Graphs work for two variables; level curves and level surfaces work in general.
  \item Open sets give room around each point, so we can move in every direction.
  \item Partial derivatives are one-variable derivatives along coordinate directions.
  \item The gradient collects them: $\grad f=(D_1f,\dots,D_nf)$.
  \item Differentiable means $f(X+H)-f(X)=\grad f(X)\cdot H+o(\norm H)$.
  \item Continuous partial derivatives imply differentiability; existing partials alone do not.
\end{itemize}
\end{keyidea}

\section{Exam-style questions}

\begin{exercise}[1]
Find and sketch the domain of $f(x,y)=\dfrac{\sqrt{9-x^2-y^2}}{\ln(x^2+y^2)}$.
\end{exercise}

\begin{exercise}[1]
Describe the level curves of $f(x,y)=y-x^2$ for $c=-1,0,1,2$ and sketch the contour map.
What does the graph look like?
\end{exercise}

\begin{exercise}[1]
Compute $\grad f$ for $f(x,y,z)=x e^{yz}+z\ln(1+x^2)$, then evaluate it at $(1,0,2)$.
\end{exercise}

\begin{exercise}[2]
Let $f(x,y)=x^3-3xy^2$. Show directly from the definition (expanding $f(X+H)-f(X)$) that $f$
is differentiable at $X=(1,1)$, identifying the linear part and checking the error is
$o(\norm H)$.
\end{exercise}

\begin{exercise}[2]
Is the set $U=\{(x,y): 1<x^2+y^2<4\}$ open in $\R^2$? Justify by giving, for a point
$P\in U$ with $\norm P=b$, an explicit radius $c$ such that the ball of radius $c$ around $P$
lies in $U$.
\end{exercise}

\begin{exercise}[2]
Let $f(x,y)=\dfrac{x^2y}{x^4+y^2}$ for $(x,y)\neq(0,0)$ and $f(0,0)=0$.
\begin{enumerate}[label=(\alph*)]
  \item Compute $D_1f(0,0)$ and $D_2f(0,0)$ from the definition.
  \item Show that $f\to0$ along every straight line through the origin.
  \item Show that $f$ is nevertheless not continuous at $O$ (try the parabola $y=x^2$).
\end{enumerate}
\end{exercise}

\begin{exercise}[3]
Let $f(x,y)=\sqrt{\abs{xy}}$. Show that $D_1f(0,0)=D_2f(0,0)=0$ but that $f$ is not
differentiable at the origin. (Hint: if it were, the remainder $f(h,k)$ would have to be
$o(\norm H)$. Test $h=k$.)
\end{exercise}

\begin{exercise}[3]
Let $f(x,y)=(x^2+y^2)\sin\!\dfrac{1}{\sqrt{x^2+y^2}}$ for $(x,y)\neq O$ and $f(O)=0$.
Show that $f$ is differentiable at $O$ with $\grad f(O)=O$, although $D_1f$ is not
continuous at $O$. What does this say about the converse of Theorem~\ref{thm:C1}
\end{exercise}
